% ======================================================================
%  MineGuard: IEEE 2-Page Conference Paper
%  Compile: pdflatex mineguard_ieee_report.tex (run twice)
% ======================================================================
\documentclass[conference]{IEEEtran}

\usepackage[utf8]{inputenc}
\usepackage[T1]{fontenc}
\usepackage{amsmath}
\usepackage{graphicx}
\usepackage{cite}
\usepackage{booktabs}
\usepackage{array}
\usepackage[hidelinks]{hyperref}

\begin{document}

\title{MineGuard: A Web-Based InSAR System Focused on
Horizontal East-West Mine Deformation Monitoring}

\author{
  \IEEEauthorblockN{Ujjwal Rawat}
  \IEEEauthorblockA{
    Dept.\ of Electronics \& Communication Engineering\\
    NIT Rourkela — Roll No.\ 225EC6010
  }
}

\maketitle

% ── Abstract ─────────────────────────────────────────────────────────
\begin{abstract}
Ground subsidence at active mine sites poses severe safety risks, but
it is the horizontal lateral sliding that typically serves as the
earliest precursor to catastrophic slope failure. Because conventional
sensors and standard Line-of-Sight (LOS) radar primarily measure vertical
settlement, these critical early warnings are often missed. This paper
presents \textbf{MineGuard}, a full-stack web application that
automates the Interferometric SAR (InSAR) pipeline using free
Sentinel-1 data, with a specific focus on isolating true East-West
horizontal deformation. By implementing a two-orbit Ascending (ASC)
and Descending (DSC) geometric decomposition, the system strips away
vertical bias to reveal pure lateral ground drift. Applied to the
Jharia Coalfield (India) over a 2017–2022 period, the application
successfully identified critical lateral landslide mechanisms,
detecting extreme eastward shifts up to $+182.2$\,mm at specific 
hotspots that would have been mischaracterized by vertical-only monitoring.
\end{abstract}

\begin{IEEEkeywords}
InSAR, Sentinel-1, horizontal deformation, lateral sliding, 2D decomposition.
\end{IEEEkeywords}

% ── I. Introduction ──────────────────────────────────────────────────
\section{Introduction}

Mining excavations progressively destabilise surrounding rock masses,
triggering both downward subsidence and lateral spreading. In geotechnical
engineering, horizontal (East-West) displacement along failure planes
is consistently recognized as the earliest precursor to slope collapse.
Once lateral sliding initiates, intervention windows are short. Yet, this
horizontal movement remains largely invisible to conventional vertical
sensors and standard single-pass radar configurations \cite{cascini2010}.

Spaceborne InSAR with Sentinel-1 offers synoptic, millimetre-scale
deformation maps at no data acquisition cost \cite{massonnet1998,
berardino2002}. However, standard InSAR workflows yield a 1-dimensional
Line-of-Sight (LOS) measurement heavily biased toward vertical motion.
Extracting true horizontal drift requires complex multi-orbit algebra
and specialist expertise unavailable to on-site engineers.
MineGuard eliminates this barrier. By embedding algorithmic 2D vector
decomposition inside a guided, automated web application, MineGuard
shifts the monitoring focus from simple subsidence to critical lateral
failure detection, providing engineers with automated, directionally-accurate
geotechnical early warnings.

% ── II. Methodology ──────────────────────────────────────────────────
\section{Methodology}

\subsection{Data Acquisition and Pre-Processing}
Sentinel-1 IW-SLC scenes are queried from the Copernicus Data Space
Ecosystem (CDSE) API. To enable horizontal decomposition, both Ascending
(ASC) and Descending (DSC) orbit passes are retrieved and assembled
into pairs with temporal baselines $\Delta t \leq 12$\,days.

Each pair is automatically processed through an ESA SNAP GPT pipeline:
orbit correction, back-geocoding, Enhanced Spectral Diversity (ESD),
interferogram formation, flat-earth/DEM phase removal, Goldstein
adaptive phase filtering, Snaphu phase unwrapping, and terrain
correction. Unwrapping removes the $\pm\pi$ wrapped-phase
ambiguity, enabling robust measurement of cumulative deformations.

\subsection{2D East-West Vector Decomposition}
A single satellite pass linearly couples horizontal and vertical
motion into one LOS orientation. Because Sentinel-1 occupies a
near-polar orbit, it is highly sensitive to East-West and vertical
movement, but largely blind to North-South motion. Therefore, the
primary objective of geometric decomposition is to mathematically
isolate the East-West drift ($d_E$) from the vertical settlement ($d_U$).

Combining matched ASC and DSC pairs forms a $2\times2$ linear system
per epoch \cite{wright2004}:
\begin{equation}
  d_{\mathrm{LOS}} = -\sin\theta\cos\alpha\cdot d_E
                   + \cos\theta\cdot d_U
  \label{eq:obs}
\end{equation}
where $\theta$ is the incidence angle ($40^\circ$ ASC, $35^\circ$
DSC), and $\alpha$ is the satellite heading ($350^\circ$ ASC,
$190^\circ$ DSC). Solving this system via Cramer's Rule yields
the pure lateral East-West drift vector, untainted by vertical
compaction artifacts.

\subsection{Horizontal-Focused Risk Classification}
While traditional systems trigger alerts based on total subsidence,
MineGuard's automated risk engine heavily weights the horizontal
magnitude. Cumulative E-W displacement defines the site-specific
lateral risk index (LOW $<30$\,mm, MEDIUM $30$--$80$\,mm, HIGH
$80$--$150$\,mm, CRITICAL $\geq150$\,mm).

% ── III. Results ─────────────────────────────────────────────────────
\section{Experimental Results}

The pipeline was deployed on the Jharia Coalfield, India
($23.732^\circ$\,N, $86.438^\circ$\,E)—a region notorious for intense
underground fires and severe surface instability. Monitoring focused
on tracking the lateral structural integrity of hotspots P1 through
P5 spanning 2017 to 2022. Table~\ref{tab:results} and Fig.~\ref{fig:plots}
detail the horizontal drift dynamics.

\begin{table}[h]
  \caption{Cumulative Horizontal Deformation Focus (Jharia, 2017–2022)}
  \label{tab:results}
  \centering
  \setlength{\tabcolsep}{6pt}
  \begin{tabular}{lccl}
    \toprule
    \textbf{Hotspot} & \textbf{East-West Drift} & \textbf{Vertical Subs.} & \textbf{Lateral Risk} \\
    \midrule
    P1 & $+315.2$\,mm & $-552.8$\,mm & CRITICAL \\
    P3 & $-188.0$\,mm & $-537.8$\,mm & CRITICAL \\
    P5 & $-178.9$\,mm & $-895.4$\,mm & CRITICAL \\
    P2 & $-22.4$\,mm  & $-397.4$\,mm & LOW \\
    P4 & $+14.5$\,mm  & $-211.0$\,mm & LOW \\
    \bottomrule
  \end{tabular}
\end{table}

\begin{figure}[h]
  \centering
  \includegraphics[width=0.7\linewidth]{MineGuard_Proto/results/hotspot_map.png}
  
  \vspace{1.5mm}
  \includegraphics[width=0.7\linewidth]{MineGuard_Proto/results/plot_01_horizontal_displacement.png}
  
  \vspace{1.5mm}
  \includegraphics[width=0.7\linewidth]{MineGuard_Proto/results/plot_03_3d_trajectory.png}
  \caption{Top: Satellite basemap showing selected hotspots. Middle: Cumulative East-West horizontal displacement isolating lateral shift. Bottom: 3D vector displacement trajectory mapping concurrent vertical and lateral motion over time.}
  \label{fig:plots}
\end{figure}

Shifting the analytical focus from vertical to horizontal completely
alters the geotechnical threat assessment. The decomposition reveals
severe mechanical diversity across the open pit. While all points
experienced massive vertical collapse (e.g., P5 subsiding an extreme
$-895.4$\,mm), the horizontal tracking captures distinct lateral
landslide mechanisms sliding in opposite directions.

P1 exhibits a massive and dangerous $+315.2$\,mm eastward shift,
representing an active lateral separation block. Conversely, P3
and P5 are accelerating westward ($-188.0$\,mm and $-178.9$\,mm
respectively). From a slope stability perspective, this severe
horizontal shearing (the pit walls spreading apart) represents an
imminent catastrophic landslide threat that could trigger sudden failure,
yet would be heavily misunderstood if observing only vertical subsidence. 
Furthermore, while P4 shows some vertical settlement, its lateral shift 
($+14.5$\,mm) is minimal, signifying centralized compaction rather than 
lateral shedding. By isolating the $d_E$ vector, the system correctly 
pinpoints the divergent lateral blocks.

% ── IV. Conclusion ───────────────────────────────────────────────────
\section{Conclusion}

MineGuard demonstrates that critical horizontal mine deformation
data can be extracted and delivered automatically through a robust,
user-friendly web interface. By implementing a 2D ASC+DSC vector
decomposition, the system successfully strips away vertical biases
to isolate pure East-West lateral sliding. Demonstrated over the
Jharia Coalfield, this focus correctly identified severe divergent
landslide threats ($+315.2$\,mm eastward at P1 and $-188.0$\,mm
westward at P3) that conventional vertical monitoring would have
obscured. This capability empowers field engineers to detect slope
failure precursors earlier, vastly improving mine-site safety
interventions.

% ── References ───────────────────────────────────────────────────────
\begin{thebibliography}{1}

\bibitem{massonnet1998}
D.~Massonnet and K.~L.~Feigl, ``Radar interferometry and its
application to changes in the earth's surface,''
\textit{Rev.\ Geophys.}, vol.~36, no.~4, 1998.

\bibitem{berardino2002}
P.~Berardino et al., ``A new algorithm for surface deformation
monitoring based on small baseline differential SAR interferograms,''
\textit{IEEE Trans.\ Geosci.\ Remote Sens.}, vol.~40, no.~11, 2002.

\bibitem{wright2004}
T.~J.~Wright, B.~E.~Parsons, and Z.~Lu, ``Toward mapping surface
deformation in three dimensions using InSAR,''
\textit{Geophys.\ Res.\ Lett.}, vol.~31, no.~1, 2004.

\bibitem{cascini2010}
L.~Cascini et al., ``Advanced low- and full-resolution DInSAR map
generation for slow-moving landslide analysis,''
\textit{Eng.\ Geol.}, vol.~112, 2010.

\end{thebibliography}

\end{document}
